\documentclass[a4paper,10pt]{article}
\usepackage[utf8]{inputenc}
\usepackage[margin=1in]{geometry}
\usepackage{tcolorbox}
\usepackage{lmodern}
\renewcommand{\familydefault}{\sfdefault}
\usepackage{graphicx}
\usepackage{parskip}

% Configure the style for the cards
\tcbset{
  colback=white,
  colframe=black,
  width=7cm, % Set square card width
  height=7cm, % Set square card height
  halign=center, % Center align text
  valign=center, % Center align text vertically
  fonttitle=\bfseries\fontsize{14pt}{16pt}\selectfont,
  fontupper=\fontsize{12pt}{14pt}\selectfont,
  coltitle=white, % Title text color
}

\begin{document}

\section*{Categorie 1: Onderdelen}

% Onderdelen (blauw titelbalk)
\tcbset{colbacktitle=blue!70}
\begin{center}
\begin{tabular}{c@{\hskip 1cm}c}
  \begin{tcolorbox}[title=ALU \\ (Arithmetic Logic Unit)]
    Voert berekeningen en logische operaties uit, zoals optellen (ADD) of vergelijken.
  \end{tcolorbox} &
  \begin{tcolorbox}[title=CU \\ (Control Unit)]
    Coördineert de werking van de CPU door instructies te decoderen en onderdelen aan te sturen.
  \end{tcolorbox} \\[1cm]
  \begin{tcolorbox}[title=Registers]
    Slaan gegevens tijdelijk op die snel toegankelijk moeten zijn voor de CPU.
  \end{tcolorbox} &
  \begin{tcolorbox}[title=Geheugen (RAM)]
    Slaat programma-instructies en gegevens op.
  \end{tcolorbox} \\[1cm]
  \begin{tcolorbox}[title=Gegevensbus]
    Verplaatst gegevens tussen onderdelen.
  \end{tcolorbox} &
  \begin{tcolorbox}[title=Adresbus]
    Geeft aan waar in het geheugen gegevens of instructies te vinden zijn.
  \end{tcolorbox}
\end{tabular}
\end{center}

\newpage

\section*{Categorie 2: Acties}

% Acties (groen titelbalk)
\tcbset{colbacktitle=red!25!green!70!blue!80}
\begin{center}
\begin{tabular}{c@{\hskip 1cm}c}
  \begin{tcolorbox}[title=Fetch]
    Haalt een instructie op uit het geheugen via de adresbus.
  \end{tcolorbox} &
  \begin{tcolorbox}[title=Decode]
    De CU vertaalt de opgehaalde instructie in acties.
  \end{tcolorbox} \\[1cm]
  \begin{tcolorbox}[title=Execute]
    De ALU voert de bewerking van de instructie uit (bijvoorbeeld optellen of vergelijken).
  \end{tcolorbox} &
  \begin{tcolorbox}[title=Store]
    Het resultaat van de bewerking wordt opgeslagen in een register of in het geheugen.
  \end{tcolorbox} \\[1cm]
  \begin{tcolorbox}[title=Load]
    Gegevens worden uit het geheugen geladen naar een register.
  \end{tcolorbox} &
  \begin{tcolorbox}[title=Branch]
    De CU bepaalt of de uitvoering naar een andere instructie moet springen op basis van een voorwaarde.
  \end{tcolorbox}
\end{tabular}
\end{center}

\newpage

\section*{Categorie 3: Gegevens en Instructies}

% Gegevens en Instructies (oranje titelbalk)
\tcbset{colbacktitle=orange!70}
\begin{center}
\begin{tabular}{c@{\hskip 1cm}c}
  \begin{tcolorbox}[title=Instructie]
    ADD R1, R2 → Optel de waarde van register R1 bij R2.
  \end{tcolorbox} &
  \begin{tcolorbox}[title=Gegevens]
    42 → Een getalwaarde die verwerkt moet worden.
  \end{tcolorbox} \\[1cm]
  \begin{tcolorbox}[title=Adres]
    0x002A → De locatie in het geheugen waar de instructie of gegevens staan.
  \end{tcolorbox} &
  \begin{tcolorbox}[title=Resultaat]
    84 → Het resultaat van een berekening of operatie.
  \end{tcolorbox}
\end{tabular}
\end{center}

\end{document}